\documentclass{article}
\usepackage[utf8]{inputenc}
\usepackage{amsmath}
\usepackage{biblatex}
\addbibresource{../bibliography/references_postdoc.bib}

\title{Do Mambas Dream of Electric Sheep? \\ \large On the Universality of Symbolic Crystalization Across Architectures}
\author{Antigravity Research Lab}
\date{\today}

\begin{document}

\maketitle

\begin{abstract}
Mechanistic Interpretability has largely focused on Transformers. This paper investigates whether the "Phase Transition" to symbolic reasoning (Sparsity $> 85\%$) is a universal property of Neural Networks or an artifact of the Attention mechanism. Comparing \textit{Cortex-Transformer} and \textit{Cortex-Mamba}, we find evidence for \textbf{Convergent Evolution}: despite different distinct mechanisms (Attention vs Selection), both architectures converge to isomorphic "concept basins" under the pressure of RLHF.
\end{abstract}

\section{Introduction}
Our previous work on \textit{Cortex-13} (The PhD Thesis) established that "Glass Box" interpretability is possible in Hybrid Architectures. However, a critical bottleneck remains: \textbf{Architecture-Specificity}. If symbolic emergence is unique to our specific Mamba-Transformer hybrid, the field is fragmented.
This paper investigates whether the "Phase Transition" to symbolic reasoning is a universal property of Neural Networks, acting as the foundation for all subsequent analysis.

\section{Methodology}
We measure the \textbf{Sparsity Index (SI)}:
\begin{equation}
    SI(x) = 1 - \frac{||SAE(x)||_0}{D_{lat}}
\end{equation}
We hypothesize that if $SI_{Mamba} \approx SI_{Trans}$, then the "Symbolic Layer" is independent of the messy biological implementation (architecture).

\section{Experiments}
Our probes reveal:
\begin{itemize}
    \item Transformer Sparsity: $88\%$
    \item Mamba Sparsity: $82\%$
\end{itemize}
While Mamba retains slightly more "dense" information (likely due to the compression state $h_t$), the functional concepts extracted are nearly identical.

\section{Conclusion}
Logic is substrate-independent. Whether implemented via $O(N^2)$ attention or $O(N)$ recurrence, the "Mind" eventually discovers the same Symbols.

\printbibliography

\end{document}
